\documentclass[runningheads]{llncs}

\usepackage[T1]{fontenc}
% T1 fonts will be used to generate the final print and online PDFs,
% so please use T1 fonts in your manuscript whenever possible.
% Other font encondings may result in incorrect characters.

%% \usepackage{graphicx}
%% % Used for displaying a sample figure. If possible, figure files should
%% % be included in EPS format.

\usepackage{hyperref}
\usepackage[nameinlink]{cleveref}
% If you use the hyperref package, please uncomment the following two lines
% to display URLs in blue roman font according to Springer's eBook style:
\usepackage{color}
\renewcommand\UrlFont{\color{blue}\rmfamily}
\urlstyle{rm}

\title{Finite Functional Logic Programming}
% If the paper title is too long for the running head, you can set
% an abbreviated paper title here
%\titlerunning{Abbreviated paper title}

\author{Michael Arntzenius\orcidID{0009-0002-0417-5636}%%  \and
%% Second Author\inst{2,3}\orcidID{1111-2222-3333-4444} \and
%% Third Author\inst{3}\orcidID{2222--3333-4444-5555}
}
%
\authorrunning{M. Arntzenius}
% First names are abbreviated in the running head.
% If there are more than two authors, 'et al.' is used.
%
\institute{UC Berkeley, Berkeley CA, USA\\
\email{daekharel@gmail.com}}


\begin{document}
\maketitle              % typeset the header of the contribution

\begin{abstract}
  I propose unifying functional and logic/relational programming by treating predicates/relations as functions equipped with their support -- the set of inputs whose output is non-false.
  For instance, the logic programming language Datalog, in which all relations are finite, is a language of finitely supported boolean functions.
  %
  Finite support allows representing the function, not as code, but as data -- as a table of input-output pairs -- making bulk operations more efficient.
%
  If we combine finitely supported functions, represented as data, and higher order functions, represented as code, in the same language, we get what I call \emph{finite functional programming.}
  %
  To this end, I demonstrate a simple type system to check finite support by treating variable grounding as a graded effect.
  %Aggregation and other Datalog extensions (such as bag and semiring semantics) are naturally handled by generalizing from boolean functions to arbitrary pointed output types.
  %% \todo{do I actually do this? $\rightarrow$}
  %% I sketch future work on semantics and implementation strategies, particularly for recursion.

\keywords{First keyword  \and Second keyword \and Another keyword.}
\end{abstract}

\section{First Section}


%% END MATTER
\begin{credits}
\subsubsection{\ackname} A bold run-in heading in small font size at the end of the paper is
used for general acknowledgments, for example: This study was funded
by X (grant number Y).

\subsubsection{\discintname}
It is now necessary to declare any competing interests or to specifically
state that the authors have no competing interests. Please place the
statement with a bold run-in heading in small font size beneath the
(optional) acknowledgments\footnote{If EquinOCS, our proceedings submission
system, is used, then the disclaimer can be provided directly in the system.},
for example: The authors have no competing interests to declare that are
relevant to the content of this article. Or: Author A has received research
grants from Company W. Author B has received a speaker honorarium from
Company X and owns stock in Company Y. Author C is a member of committee Z.
\end{credits}

%% %
%% % ---- Bibliography ----
%% %
%% % BibTeX users should specify bibliography style 'splncs04'.
%% % References will then be sorted and formatted in the correct style.
%% %
%% % \bibliographystyle{splncs04}
%% % \bibliography{mybibliography}
%% %
%% \begin{thebibliography}{8}
%% \bibitem{ref_article1}
%% Author, F.: Article title. Journal \textbf{2}(5), 99--110 (2016)

%% \bibitem{ref_lncs1}
%% Author, F., Author, S.: Title of a proceedings paper. In: Editor,
%% F., Editor, S. (eds.) CONFERENCE 2016, LNCS, vol. 9999, pp. 1--13.
%% Springer, Heidelberg (2016). \doi{10.10007/1234567890}

%% \bibitem{ref_book1}
%% Author, F., Author, S., Author, T.: Book title. 2nd edn. Publisher,
%% Location (1999)

%% \bibitem{ref_proc1}
%% Author, A.-B.: Contribution title. In: 9th International Proceedings
%% on Proceedings, pp. 1--2. Publisher, Location (2010)

%% \bibitem{ref_url1}
%% LNCS Homepage, \url{http://www.springer.com/lncs}, last accessed 2023/10/25
%% \end{thebibliography}

\end{document}
